% Generated draft: requires mathematical and source review.
\documentclass[11pt]{article}
\usepackage[margin=25mm]{geometry}
\usepackage{fontspec}
\setmainfont{Latin Modern Roman}
\usepackage{amsmath,amssymb,mathtools,unicode-math}
\setmathfont{Latin Modern Math}
\usepackage{enumitem,needspace,xurl,hyperref}
\setlength{\parindent}{0pt}
\setlength{\parskip}{5pt}
\setlength{\emergencystretch}{3em}
\newcommand{\N}{\mathbb{N}}
\newcommand{\Z}{\mathbb{Z}}
\newcommand{\Q}{\mathbb{Q}}
\newcommand{\R}{\mathbb{R}}
\newcommand{\C}{\mathbb{C}}
\newcommand{\F}{\mathbb{F}}
\newcommand{\A}{\mathbb{A}}
\newcommand{\PP}{\mathbb{P}}
\newcommand{\E}{\mathbb{E}}
\newcommand{\eps}{\varepsilon}
\newcommand{\dd}{\,\mathrm{d}}
\begin{document}
\section*{UnsolvedMath ID: 30004728}
\textbf{Source problem code:} catalogue entry
\subsection{Definitions and mathematical statement}
A closed manifold is compact and has no boundary. To interpret Riemannian metrics, take the manifolds to be smooth; the supplied statement does not explicitly specify the smooth category. For a smooth manifold \(X\), define
\[
R^+(X)=\{g\in\Gamma^\infty(\operatorname{Sym}^2T^*X): g\text{ is positive definite and }\operatorname{scal}_g(x)>0\text{ for every }x\in X\},
\]
with the \(C^\infty\) topology. Here \(\operatorname{scal}_g\) denotes scalar curvature. This is the space of metrics themselves, without taking a quotient by diffeomorphisms. The topology is a standard interpretive convention and is not specified in the supplied record.

Under the usual stable normal convention, let \(\nu_X:X\to BO\) classify the stable normal bundle of \(X\), where \(BO\) is the stable classifying space for real vector bundles. A normal \(2\)-type of \(X\) is represented by a fibration \(p:B\to BO\) and a lift \(\bar\nu_X:X\to B\), with \(p\circ\bar\nu_X\simeq\nu_X\), such that \(\bar\nu_X\) is \(3\)-connected and each homotopy fiber of \(p\) has vanishing homotopy groups in degrees at least \(3\). A \(3\)-connected map induces a bijection on path components, isomorphisms on \(\pi_1\) and \(\pi_2\), and a surjection on \(\pi_3\), at every basepoint. Having the same normal \(2\)-type means that the corresponding fibrations over \(BO\) are fiber homotopy equivalent; equivalently, a common such fibration can be used for both normal lifts. The source does not specify its convention for normal \(2\)-types, connectedness, or additional tangential structures; this definition records the usual unstructured stable normal interpretation, subject to source review.

The supplied conjecture is the following: for every integer \(d\geq5\) and every pair of closed smooth \(d\)-manifolds \(M,N\) with the same normal \(2\)-type,
\[
R^+(M)\neq\varnothing\quad\text{and}\quad R^+(N)\neq\varnothing
\quad\Longrightarrow\quad
R^+(M)\simeq R^+(N).
\]
Explicitly, the conclusion asks for continuous maps \(F:R^+(M)\to R^+(N)\) and \(G:R^+(N)\to R^+(M)\) satisfying
\[
G\circ F\simeq\operatorname{id}_{R^+(M)},\qquad
F\circ G\simeq\operatorname{id}_{R^+(N)}.
\]
No connectedness, orientability, spin structure, or normal bordism relation between \(M\) and \(N\) is explicitly assumed in the supplied formulation. Whether the original discussion imposes any such qualifications remains unresolved. The source labels this assertion a conjecture and supplies an open-status assessment; neither constitutes an independently verified mathematical conclusion here.
\subsection{Short English statement}
The supplied conjecture asserts that, in every dimension \(d\geq5\), two closed manifolds with the same normal \(2\)-type have homotopy equivalent spaces of positive scalar curvature metrics, provided both spaces are nonempty. Smoothness, the \(C^\infty\) topology, and the precise normal \(2\)-type convention require confirmation against the original source.
\subsection{Sources}
\begin{itemize}
\item Supplied local UnsolvedMath registry record, permanent ID \(30004728\), title \emph{Homotopy Invariance of Positive-Scalar-Curvature Metric Spaces}, source problem code \texttt{catalogue entry}, in the Oberwolfach Reports Open Problems collection. The record gives proposed year \(2021\) and names no individual proposer.
\item Supplied report citation: \emph{Analysis, Geometry and Topology of Positive Scalar Curvature Metrics} (2022), Oberwolfach Reports, pp.\ 1579--1649, DOI \texttt{10.4171/owr/2021/30}. \url{https://doi.org/10.4171/owr/2021/30}. The record describes Oberwolfach Reports as documenting mathematical workshops and their open-problem discussions. The cited report has not been independently consulted here.
\item Supplied literature source: \emph{Journal of the European Mathematical Society, positive scalar curvature spaces}. \url{https://ems.press/content/serial-article-files/31422}.
\item Supplied literature source: \emph{Oberwolfach Reports 2022}. \url{https://ems.press/content/serial-article-files/46909}.
\item Supplied literature assessment dated exactly \texttt{2026-08-22}: the same-normal-two-type homotopy-equivalence conjecture remains open, with existing results covering selected normal two-types and special fundamental groups rather than the stated generality. This is supplied dated curation derived from the OWR record, not an independently verified assessment; the record explicitly requests independent verification of time-sensitive status and cited claims.
\item The registry labels the statement \texttt{corrected\_verified}, explaining that a publication-format audit preserved the mathematical content while repairing notation or typography. This supplied audit label does not establish mathematical correctness or verify the literature status.
\end{itemize}
\end{document}
